上一单元把表示交给了优化，接下来要说清楚优化到底在追什么，以及追的力气怎样分配到每一个参数上。结构规定网络能输出什么，损失函数规定哪些输出值得寻找，反向传播计算参数的改变会如何影响损失——三件事各管一段，混起来谈就查不出问题出在哪里。

第一篇处理损失。平方损失的最优预测是条件均值，绝对损失落在条件中位数，分位数损失落在指定分位点；换个角度看，这些损失分别是高斯、Laplace 和非对称 Laplace 噪声假设的负对数，交叉熵则对应 Bernoulli 或 Categorical 分布。换损失就是换假设，而网络自己不会提醒你假设已经被数据推翻。第二篇处理梯度。反向传播就是链式法则按从右往左的顺序乘，一趟前向加一趟反向拿到全部梯度，总代价与一次前向同阶。矩阵微积分那一部分已经给出一般计算图上的推导和代价结构，这里补的是另一半：一次完整的数值走查，以及实践中怎样定位梯度问题。

梯度检查只提供局部数值证据。通过检查不代表损失语义正确，也不代表训练稳定；检查失败时的排查顺序是形状、广播带来的隐式求和、批规约因子、不可微点、有限差分步长与浮点精度。

\subsection{怎么读这一章}

两篇都是主线。若眼下只想排查实现错误，可以先读第二篇；但在解释训练结果之前必须回到第一篇，因为梯度算对只说明忠实优化了当前这个损失，不说明这个损失对应你真正要的目标。

\subsection{本章篇目}

以下篇目按概念依赖排列；若从具体问题进入，也可以先打开最接近当前困惑的一篇，再沿篇内链接回补。

\begin{itemize}
\tightlist
\item \hyperref[sec:14-Deep-Learning/02-Losses-and-Backpropagation/01]{``损失函数决定网络在学习什么''}；
\item \hyperref[sec:14-Deep-Learning/02-Losses-and-Backpropagation/02]{``一次完整的前向与反向传播''}。
\end{itemize}

目标定准、梯度算对之后，训练仍可能失败，那部分归下一单元。
